\scanpage{173}
% archive foliation: 172
\uline{Le foncteur de Picard}
\margin{Faits élémentaires sur les classes de diviseurs, faisceaux inversibles, leurs propriétés etc}

Soit $X$ un $S$-préschéma, on pose
\[
\mathrm{Prépic}(X/S)=\deleted{H^1(X,\underline O_X^*)}\ \operatorname{Pic}(X\deleted{\unsure{/S}})=H^1(X,\underline O_X^*)
\]
\[
\underline{\underline{\mathrm{Prépic}}}_{X/S}(S')=\mathrm{Prépic}(X'/S')=H^1(X',\underline O_{X'}^*)
\quad\text{où }X'=X\times_S S'
\]
donc $\underline{\underline{\mathrm{Prépic}}}_{X/S}$ est bien un foncteur \uline{contravariant} sur les préschémas $S'$
\[
\mathrm{Prépic}'(X/S)=\Gamma(S,\underline H^1_{X/S}(\underline O_X^*))
\]
\[
\underline{\underline{\mathrm{Prépic}}}'_{X/S}(S')=\Gamma(S',\underline H^1_{X'/S'}(\underline O_{X'}^*))
\quad\text{où }X'=X\times_S S'
\]
On passe de $\underline{\underline{\mathrm{Prépic}}}$ à $\underline{\underline{\mathrm{Prépic}}}'$ en effectuant la descente par recouvrements ouverts.
\margin{expliciter les \unsure{objets}, \unsure{l'épimorphisme} $\operatorname{Pic}(X\deleted{/S})$ défini par $\operatorname{Pic}(X)\twoheadrightarrow\mathrm{Prépic}'(X/S)$}

\uline{Proposition 1} Supposons $\underline H^0_{X/S}(\underline O_X)=\underline O_S$. Alors on a un \deleted{isomorphisme}\inserted{monomorphisme} canonique
\[
\mathrm{Prépic}'(X/S)\longleftarrow\mathrm{Prépic}(X/S)/\operatorname{Pic}\mathrm{abs}(S)
\]
qui est un isomorphisme si $X/S$ a une section.

Résulte de l'assertion plus précise
\margin{mettre ce \deleted{corollaire} sous forme d'une suite exacte plutôt}

\uline{Corollaire 1} Le foncteur $\underline F\to\underline F\otimes_S X$ des faisceaux loc. libres de $S$ dans les faisceaux loc. libres de $X$ \deleted{définit} est \uline{pleinement} fidèle, l'image essentielle est formée des faisceaux loc. libres sur $X$ qui sont triviaux loc. en dessus de $\unsure{Y|S}$, le foncteur inverse sur ce domaine est le foncteur $\underline H^0_{X/S}$.
\margin{\unsure{val. pour une cond. plus forte} \illegible{2} \unsure{projection pleinement}}
\margin{\unsure{dire que le} \illegible{2} \unsure{est un} \illegible{1}}

